% GENERATED FILE. Edit canonical lessons, then run npm run generate:books.
% canonical-source-hash: fnv1a64:79492f8579ad4ebd
% canonical-lessons: TE-C114-galipatam, TE-C114-cokka, TE-C114-cira, TE-C114-pance, TE-C114-ceppulu

\chapter{Kite, Shirt, Sari, Dhoti, Sandals}
\label{ch:te-kite-shirt-sari-dhoti-sandals}
\hlchaptermodality{114}

\begin{chapteropening}
\textbf{By the end of this chapter:} \emph{I can say a kite, a shirt, a sari, a dhoti and sandals.}

Complete the last lesson of chapter 114: i can say a kite, a shirt, a sari, a dhoti and sandals.
\end{chapteropening}

\section[\texorpdfstring{\te{గాలిపటం} (gālipaṭaṁ)}{గాలిపటం (gālipaṭaṁ)}]{\te{గాలిపటం} (gālipaṭaṁ) — a kite}

\label{lesson:TE-C114-galipatam}



\begin{quote}
Before the new one: say the Telugu for a doll; a picture, then the Telugu for a ball.
\end{quote}

\subsection*{You'll want to know: \te{గాలిపటం}}
\textbf{\te{గాలిపటం}} — \emph{gālipaṭaṁ} — \textquotedblleft{}a kite\textquotedblright{}.

\begin{grammarlens}[title={What you've built}]
One more word for this chapter.
\end{grammarlens}

\subsection*{Your turn}
\begin{itemize}
\raggedright
  \item \emph{Say it:} \emph{gālipaṭaṁ}
  \item \emph{Say it:} \emph{gālipaṭaṁ}, once more
  \item \emph{Say it:} say \emph{gālipaṭaṁ}
  \item \emph{Recall:} say the Telugu for a doll; a picture, then the Telugu for a ball, then say \emph{gālipaṭaṁ} again
\end{itemize}

\begin{tcolorbox}[breakable,colback=teal!4,colframe=teal!35!black,title={Before you move on}]
What does \te{గాలిపటం} mean? (\textquotedblleft{}A kite\textquotedblright{}.) Say it once more.
\end{tcolorbox}
\section[\texorpdfstring{\te{చొక్కా} (cokkā)}{చొక్కా (cokkā)}]{\te{చొక్కా} (cokkā) — a shirt}

\label{lesson:TE-C114-cokka}



\begin{quote}
Before the new one: say the Telugu for a ball, then the Telugu for a kite.
\end{quote}

\subsection*{You'll want to know: \te{చొక్కా}}
\textbf{\te{చొక్కా}} — \emph{cokkā} — \textquotedblleft{}a shirt\textquotedblright{}.

\begin{grammarlens}[title={What you've built}]
One more word for this chapter.
\end{grammarlens}

\subsection*{Your turn}
\begin{itemize}
\raggedright
  \item \emph{Say it:} \emph{cokkā}
  \item \emph{Say it:} \emph{cokkā}, once more
  \item \emph{Say it:} say \emph{cokkā}
  \item \emph{Recall:} say the Telugu for a ball, then the Telugu for a kite, then say \emph{cokkā} again
\end{itemize}

\begin{tcolorbox}[breakable,colback=teal!4,colframe=teal!35!black,title={Before you move on}]
What does \te{చొక్కా} mean? (\textquotedblleft{}A shirt\textquotedblright{}.) Say it once more.
\end{tcolorbox}
\section[\texorpdfstring{\te{చీర} (cīra)}{చీర (cīra)}]{\te{చీర} (cīra) — a sari}

\label{lesson:TE-C114-cira}



\begin{quote}
Before the new one: say the Telugu for a kite, then the Telugu for a shirt.
\end{quote}

\subsection*{You'll want to know: \te{చీర}}
\textbf{\te{చీర}} — \emph{cīra} — \textquotedblleft{}a sari\textquotedblright{}.

\begin{grammarlens}[title={What you've built}]
One more word for this chapter.
\end{grammarlens}

\subsection*{Your turn}
\begin{itemize}
\raggedright
  \item \emph{Say it:} \emph{cīra}
  \item \emph{Say it:} \emph{cīra}, once more
  \item \emph{Say it:} say \emph{cīra}
  \item \emph{Recall:} say the Telugu for a kite, then the Telugu for a shirt, then say \emph{cīra} again
\end{itemize}

\begin{tcolorbox}[breakable,colback=teal!4,colframe=teal!35!black,title={Before you move on}]
What does \te{చీర} mean? (\textquotedblleft{}A sari\textquotedblright{}.) Say it once more.
\end{tcolorbox}
\section[\texorpdfstring{\te{పంచె} (pañce)}{పంచె (pañce)}]{\te{పంచె} (pañce) — a dhoti}

\label{lesson:TE-C114-pance}



\begin{quote}
Before the new one: say the Telugu for a shirt, then the Telugu for a sari.
\end{quote}

\subsection*{You'll want to know: \te{పంచె}}
\textbf{\te{పంచె}} — \emph{pañce} — \textquotedblleft{}a dhoti\textquotedblright{}.

\begin{grammarlens}[title={What you've built}]
One more word for this chapter.
\end{grammarlens}

\subsection*{Your turn}
\begin{itemize}
\raggedright
  \item \emph{Say it:} \emph{pañce}
  \item \emph{Say it:} \emph{pañce}, once more
  \item \emph{Say it:} say \emph{pañce}
  \item \emph{Recall:} say the Telugu for a shirt, then the Telugu for a sari, then say \emph{pañce} again
\end{itemize}

\begin{tcolorbox}[breakable,colback=teal!4,colframe=teal!35!black,title={Before you move on}]
What does \te{పంచె} mean? (\textquotedblleft{}A dhoti\textquotedblright{}.) Say it once more.
\end{tcolorbox}
\section[\texorpdfstring{\te{చెప్పులు} (ceppulu)}{చెప్పులు (ceppulu)}]{\te{చెప్పులు} (ceppulu) — sandals}

\label{lesson:TE-C114-ceppulu}



\begin{quote}
Before the new one: say the Telugu for a sari, then the Telugu for a dhoti.
\end{quote}

\subsection*{You'll want to know: \te{చెప్పులు}}
\textbf{\te{చెప్పులు}} — \emph{ceppulu} — \textquotedblleft{}sandals\textquotedblright{}.

\begin{grammarlens}[title={What you've built}]
One more word for this chapter.
\end{grammarlens}

\subsection*{Your turn}
\begin{itemize}
\raggedright
  \item \emph{Say it:} \emph{ceppulu}
  \item \emph{Say it:} \emph{ceppulu}, once more
  \item \emph{Say it:} say \emph{ceppulu}
  \item \emph{Recall:} say the Telugu for a sari, then the Telugu for a dhoti, then say \emph{ceppulu} again
\end{itemize}

\begin{tcolorbox}[breakable,colback=teal!4,colframe=teal!35!black,title={Before you move on}]
What does \te{చెప్పులు} mean? (\textquotedblleft{}Sandals\textquotedblright{}.) Say it once more.
\end{tcolorbox}
